\documentclass[a4paper,12pt]{article}

\usepackage[T2A]{fontenc}
\usepackage[utf8]{inputenc}
\usepackage[russian]{babel}

\usepackage{amsmath,amssymb,amsthm,mathtools}
\usepackage{helvet}
\renewcommand{\familydefault}{\sfdefault}
\usepackage{icomma}
\usepackage[left=18mm,right=18mm,top=18mm,bottom=20mm]{geometry}
\usepackage{microtype}
\usepackage{tikz}
\usetikzlibrary{arrows.meta,calc}
\usepackage{tkz-euclide}
\usepackage{pgfplots}
\pgfplotsset{compat=1.18}
\usepackage{tabularx,booktabs,longtable}
\usepackage{enumitem}
\usepackage{hyperref}
\usepackage{parskip}
\usepackage{fancyhdr}
\usepackage{setspace}

\definecolor{accent}{HTML}{008080}
\definecolor{darkaccent}{HTML}{004D4D}
\definecolor{textgray}{HTML}{333333}
\definecolor{lightline}{HTML}{D9EAEA}

\hypersetup{hidelinks}
\onehalfspacing
\setlength{\parindent}{0pt}
\setlength{\parskip}{6pt}

\pagestyle{fancy}
\fancyhf{}
\lhead{\textcolor{darkaccent}{Показательные неравенства}}
\rhead{\textcolor{darkaccent}{Николь}}
\cfoot{\thepage}

\newcommand{\R}{\mathbb{R}}
\newcommand{\answer}[1]{\textbf{Ответ: }#1}
\newcommand{\important}[1]{\textcolor{darkaccent}{\textbf{#1}}}

\newcommand{\sectionline}{%
  \vspace{2mm}
  \noindent\textcolor{lightline}{\rule{\linewidth}{0.8pt}}
  \vspace{1mm}
}

\begin{document}

\begin{titlepage}
\thispagestyle{empty}

\vspace*{18mm}

{\Huge\bfseries\textcolor{darkaccent}{Чек-лист}}

\vspace{8mm}

{\Large\bfseries Показательные неравенства}

\vspace{4mm}

{\large Замены, метод интервалов и аккуратное оформление решения}

\vspace{18mm}

\textbf{Ученица:} Николь

\vspace{6mm}

\textbf{Цель:} закрепить основные приёмы решения показательных неравенств и убрать типичные ошибки, которые приводят к потере баллов: неверные скобки в ответе, необоснованные числовые прямые, пропущенные ограничения и неправильное деление в неравенствах.

\vfill

{\large Лёвин Артём Александрович}

\vspace{4mm}

{\large \today}

\vspace*{15mm}

\end{titlepage}

\section*{1. Главная идея темы}

Показательное неравенство чаще всего нужно привести к одному из стандартных видов:
\[
a^x,\qquad (a^x)^2,\qquad \left(\frac ab\right)^x,\qquad t=a^x,\quad t>0.
\]

После такой замены исходное неравенство обычно превращается в линейное, квадратное или рациональное неравенство.

\sectionline

\textbf{Базовый алгоритм}

\begin{enumerate}[label=\textbf{\arabic*.}, leftmargin=8mm]
\item Перенести всё в одну сторону, если нужно получить вид
\[
F(x)>0,\qquad F(x)\ge0,\qquad F(x)<0,\qquad F(x)\le0.
\]
\item Разложить составные основания:
\[
4^x=(2^2)^x=2^{2x},\qquad
25^x=(5^2)^x=5^{2x},\qquad
10^x=(2\cdot5)^x=2^x\cdot5^x.
\]
\item Сделать степени или основания удобными для замены.
\item Выполнить замену, обязательно записав ограничение:
\[
t=a^x,\qquad t>0.
\]
\item Решить полученное неравенство методом интервалов.
\item Вернуться от \(t\) к \(x\).
\item Проверить скобки в ответе.
\end{enumerate}

\section*{2. Ошибки, которые важно не повторять}

\subsection*{Ошибка 1. Необоснованная числовая прямая}

Если в решении появляется числовая прямая с точками \(1\), \(4\), \(0\), \(2\), нужно показать, откуда эти точки взялись.

Например:
\[
\frac{4-t}{t-1}>0.
\]

Критические точки находятся так:
\[
4-t=0 \Rightarrow t=4,
\]
\[
t-1=0 \Rightarrow t=1.
\]

При этом точка \(t=1\) не входит в ответ, потому что она обращает знаменатель в ноль:
\[
t-1\ne0.
\]

\subsection*{Ошибка 2. Корень знаменателя нельзя включать в ответ}

Если точка получилась из знаменателя, она всегда выколота. Даже при нестрогом знаке \(\le\) или \(\ge\) такая точка не входит в ответ.

Например:
\[
\frac{t-4}{t-1}\ge0.
\]
Точка \(t=1\) выколота, потому что при \(t=1\) дробь не определена.

\subsection*{Ошибка 3. Нельзя заменять метод интервалов отдельными условиями}

Если после преобразований получилось
\[
x(x-2)<0,
\]
нельзя писать отдельно:
\[
x>0,\qquad x>2.
\]

Правильный ход: отметить корни \(0\) и \(2\), расставить знаки и выбрать нужный промежуток.

\[
\begin{tikzpicture}[scale=1]
\draw[-{Stealth[length=2mm]}] (-1,0)--(5,0);
\draw (1,0.1)--(1,-0.1);
\draw (3,0.1)--(3,-0.1);
\node[below] at (1,-0.1) {\(0\)};
\node[below] at (3,-0.1) {\(2\)};
\draw[fill=white,draw=black,thick] (1,0) circle (2pt);
\draw[fill=white,draw=black,thick] (3,0) circle (2pt);
\node[above] at (0,0) {\(+\)};
\node[above] at (2,0) {\(-\)};
\node[above] at (4,0) {\(+\)};
\draw[very thick,accent] (1,0.35)--(3,0.35);
\draw[fill=white,draw=accent,thick] (1,0.35) circle (2pt);
\draw[fill=white,draw=accent,thick] (3,0.35) circle (2pt);
\end{tikzpicture}
\]

Так как нужно \(x(x-2)<0\), выбираем промежуток:
\[
x\in(0;2).
\]

\subsection*{Ошибка 4. Круглая скобка вместо квадратной меняет ответ}

Если точка входит в ответ, нужна квадратная скобка:
\[
x\ge0 \quad \Longleftrightarrow \quad x\in[0;+\infty).
\]

Запись
\[
x\in(0;+\infty)
\]
означает уже другое условие:
\[
x>0.
\]

\subsection*{Ошибка 5. Корень можно извлекать из нуля}

Для корня чётной степени:
\[
\sqrt{x}
\]
область допустимых значений:
\[
x\ge0.
\]

Ноль разрешён:
\[
\sqrt0=0.
\]

Поэтому условие \(x>0\) вместо \(x\ge0\) может привести к потере точки в ответе.

\section*{3. Как оформлять систему через пересечение множеств}

Если решается система, нужно показать оба множества и их пересечение.

Например:
\[
\begin{cases}
x>1,\\
x<\log_3 4.
\end{cases}
\]

\[
\begin{tikzpicture}[scale=1]
\draw[-{Stealth[length=2mm]}] (-0.4,0)--(6,0);

\draw (1,0.1)--(1,-0.1);
\draw (4,0.1)--(4,-0.1);
\node[below] at (1,-0.1) {\(1\)};
\node[below] at (4,-0.1) {\(\log_3 4\)};

\draw[very thick,accent] (1,0.45)--(5.7,0.45);
\draw[fill=white,draw=accent,thick] (1,0.45) circle (2pt);
\node[right] at (5.7,0.45) {\(x>1\)};

\draw[very thick,darkaccent] (-0.2,0.85)--(4,0.85);
\draw[fill=white,draw=darkaccent,thick] (4,0.85) circle (2pt);
\node[left] at (-0.2,0.85) {\(x<\log_3 4\)};

\draw[very thick,black] (1,1.25)--(4,1.25);
\draw[fill=white,draw=black,thick] (1,1.25) circle (2pt);
\draw[fill=white,draw=black,thick] (4,1.25) circle (2pt);
\node[above] at (2.5,1.25) {\(\left(1;\log_3 4\right)\)};
\end{tikzpicture}
\]

\[
\answer{\displaystyle x\in\left(1;\log_3 4\right).}
\]

\section*{4. Рациональное неравенство после замены}

Рассмотрим неравенство:
\[
\frac{4-2^x}{2^x-1}>0.
\]

Делаем замену:
\[
t=2^x,\qquad t>0.
\]

Получаем:
\[
\frac{4-t}{t-1}>0.
\]

Критические точки:
\[
4-t=0 \Rightarrow t=4,
\]
\[
t-1=0 \Rightarrow t=1,\qquad t=1 \text{ не входит в область определения.}
\]

Решаем методом интервалов:
\[
\begin{tikzpicture}[scale=1]
\draw[-{Stealth[length=2mm]}] (-0.5,0)--(6,0);

\draw (1,0.1)--(1,-0.1);
\draw (4,0.1)--(4,-0.1);
\node[below] at (1,-0.1) {\(1\)};
\node[below] at (4,-0.1) {\(4\)};

\draw[fill=white,draw=black,thick] (1,0) circle (2pt);
\draw[fill=white,draw=black,thick] (4,0) circle (2pt);

\node[above] at (0.4,0) {\(-\)};
\node[above] at (2.5,0) {\(+\)};
\node[above] at (5,0) {\(-\)};

\draw[very thick,accent] (1,0.35)--(4,0.35);
\draw[fill=white,draw=accent,thick] (1,0.35) circle (2pt);
\draw[fill=white,draw=accent,thick] (4,0.35) circle (2pt);
\end{tikzpicture}
\]

Так как требуется \(>0\), получаем:
\[
t\in(1;4).
\]

Возвращаемся к \(x\):
\[
1<2^x<4.
\]

Так как
\[
1=2^0,\qquad 4=2^2,
\]
получаем:
\[
0<x<2.
\]

\[
\answer{\displaystyle x\in(0;2).}
\]

\section*{5. Когда множитель можно убрать}

Положительный множитель не влияет на знак произведения.

Например:
\[
2(4-t)>0.
\]

Так как \(2>0\), можно решать:
\[
4-t>0.
\]

То же самое верно для выражений вида:
\[
3^{x+2}>0,\qquad 5^x>0,\qquad 2^{2x}>0.
\]

Если множитель всегда положителен, его можно убрать из произведения или использовать как безопасный делитель. Если знак множителя неизвестен, так делать нельзя.

\section*{6. Деление в неравенствах}

В неравенствах нельзя бездумно делить на выражение с переменной.

\sectionline

\textbf{Можно делить}, если точно известно, что выражение положительно:
\[
3^x>0,\qquad 2^{2x}>0,\qquad (5^x)^2>0.
\]

При делении на положительное выражение знак неравенства не меняется.

\sectionline

\textbf{Нельзя делить без проверки}, если знак выражения неизвестен:
\[
x-2,\qquad x+1,\qquad 2x-5.
\]

Такие выражения могут быть положительными, отрицательными или равными нулю.

\section*{7. Аккуратная работа со степенями}

Полезные преобразования:
\[
a^{m+n}=a^m\cdot a^n,
\]
\[
a^{m-n}=\frac{a^m}{a^n},
\]
\[
(a^m)^n=a^{mn},
\]
\[
a^x b^x=(ab)^x.
\]

Например:
\[
2^{t-1}=\frac{2^t}{2},
\qquad
3^{t-4}=\frac{3^t}{3^4}.
\]

Тогда:
\[
2^{t-1}\cdot3^{t-4}
=
\frac{2^t}{2}\cdot\frac{3^t}{3^4}
=
\frac{6^t}{2\cdot3^4}.
\]

Главная проверка: если выполняется деление, нужно делить всю левую часть и всю правую часть неравенства, а не только один множитель.

\section*{8. Однородные показательные неравенства: способ 1}

Рассмотрим неравенство:
\[
5\cdot4^x-7\cdot10^x+2\cdot25^x<0.
\]

Раскладываем основания:
\[
4^x=(2^2)^x=(2^x)^2,
\]
\[
10^x=(2\cdot5)^x=2^x\cdot5^x,
\]
\[
25^x=(5^2)^x=(5^x)^2.
\]

Получаем:
\[
5(2^x)^2-7\cdot2^x\cdot5^x+2(5^x)^2<0.
\]

Делим на \((2^x)^2\). Это разрешено, потому что:
\[
(2^x)^2>0.
\]

Получаем:
\[
5-7\left(\frac52\right)^x+2\left(\frac52\right)^{2x}<0.
\]

Делаем замену:
\[
t=\left(\frac52\right)^x,\qquad t>0.
\]

Тогда:
\[
2t^2-7t+5<0.
\]

Решаем квадратное неравенство:
\[
2t^2-7t+5=0,
\]
\[
D=(-7)^2-4\cdot2\cdot5=49-40=9,
\]
\[
t_1=\frac{7-3}{4}=1,\qquad
t_2=\frac{7+3}{4}=\frac52.
\]

Так как ветви параболы направлены вверх, выражение меньше нуля между корнями:
\[
t\in\left(1;\frac52\right).
\]

Возвращаемся к \(x\):
\[
1<\left(\frac52\right)^x<\frac52.
\]

Так как
\[
1=\left(\frac52\right)^0,
\qquad
\frac52=\left(\frac52\right)^1,
\]
получаем:
\[
0<x<1.
\]

\[
\answer{\displaystyle x\in(0;1).}
\]

\section*{9. Однородные показательные неравенства: способ 2}

То же неравенство можно решить через две замены:
\[
5(2^x)^2-7\cdot2^x\cdot5^x+2(5^x)^2<0.
\]

Пусть
\[
t=2^x,\qquad k=5^x.
\]
Тогда:
\[
t>0,\qquad k>0.
\]

Получаем:
\[
5t^2-7tk+2k^2<0.
\]

Считаем это квадратным трёхчленом относительно \(t\), а \(k\) считаем коэффициентом:
\[
a=5,\qquad b=-7k,\qquad c=2k^2.
\]

Дискриминант:
\[
D=(-7k)^2-4\cdot5\cdot2k^2=49k^2-40k^2=9k^2.
\]

Корни:
\[
t_1=\frac{7k-3k}{10}=\frac{2k}{5},
\]
\[
t_2=\frac{7k+3k}{10}=k.
\]

Разложение:
\[
5t^2-7tk+2k^2
=
5\left(t-\frac{2k}{5}\right)(t-k).
\]

Возвращаемся к \(x\):
\[
5\left(2^x-\frac25\cdot5^x\right)(2^x-5^x)<0.
\]

Нули скобок:
\[
2^x=5^x
\]
и
\[
2^x=\frac25\cdot5^x.
\]

Первое уравнение:
\[
2^x=5^x.
\]
Делим на \(5^x>0\):
\[
\left(\frac25\right)^x=1.
\]
Значит:
\[
x=0.
\]

Второе уравнение:
\[
2^x=\frac25\cdot5^x.
\]
Делим на \(5^x>0\):
\[
\left(\frac25\right)^x=\frac25.
\]
Значит:
\[
x=1.
\]

Критические точки:
\[
0,\qquad 1.
\]

Проверка знаков даёт:
\[
x\in(0;1).
\]

\[
\answer{\displaystyle x\in(0;1).}
\]

\section*{10. Разложение квадратного трёхчлена}

Если квадратный трёхчлен
\[
ax^2+bx+c
\]
имеет корни \(x_1\) и \(x_2\), то:
\[
ax^2+bx+c=a(x-x_1)(x-x_2).
\]

Пример:
\[
x^2-x-6.
\]

Корни:
\[
x_1=-2,\qquad x_2=3.
\]

Тогда:
\[
x^2-x-6=(x+2)(x-3).
\]

Это полезно, потому что выражение с квадратом превращается в произведение двух простых скобок.

\section*{11. Логарифмирование}

Логарифмирование удобно, когда степени нельзя сразу привести к одному очевидному основанию.

Например:
\[
2^x=5^x.
\]

Логарифмируем обе части по основанию \(2\):
\[
\log_2 2^x=\log_2 5^x.
\]

Выносим показатель:
\[
x\log_2 2=x\log_2 5.
\]

Так как
\[
\log_2 2=1,
\]
получаем:
\[
x=x\log_2 5.
\]

Переносим всё в одну сторону:
\[
x-x\log_2 5=0,
\]
\[
x(1-\log_2 5)=0.
\]

Так как
\[
1-\log_2 5\ne0,
\]
получаем:
\[
x=0.
\]

\section*{12. Проверка перед записью ответа}

Перед финальным ответом нужно проверить:

\begin{enumerate}[label=\(\square\), leftmargin=8mm]
\item Все ли критические точки найдены.
\item Выколоты ли точки, полученные из знаменателя.
\item Записано ли условие замены \(t>0\).
\item Не потеряна ли точка \(0\) в области определения корня.
\item Верно ли выбраны круглые и квадратные скобки.
\item Выполнен ли возврат от новой переменной к \(x\).
\item Если решалась система, показано ли пересечение множеств.
\item Если было деление на выражение с переменной, объяснено ли, почему оно положительно.
\end{enumerate}

\newpage

\section*{Тренировочный блок}

\subsection*{Задание 1}

Решите неравенство:
\[
(3^x-3)(2-3^x)\ge0.
\]

\subsection*{Задание 2}

Решите неравенство:
\[
\frac{4-2^x}{2^x-1}>0.
\]

\subsection*{Задание 3}

Решите неравенство:
\[
\sqrt{x}\,(x-2)\le0.
\]

\subsection*{Задание 4}

Решите неравенство:
\[
5\cdot4^x-7\cdot10^x+2\cdot25^x<0.
\]

\subsection*{Задание 5}

Решите неравенство:
\[
7\cdot3^{2x}+84\cdot3^x7^x-3\cdot7^{2x}<0.
\]

\subsection*{Задание 6}

Решите неравенство:
\[
25^x-6\cdot5^x+5\le0.
\]

\subsection*{Задание 7}

Решите неравенство:
\[
\frac{9^x-4}{3^x-1}\ge0.
\]

\subsection*{Задание 8}

Решите неравенство:
\[
2^{x+1}+2^{x+2}-24<0.
\]

\subsection*{Задание 9}

Решите неравенство:
\[
\left(\frac13\right)^x<9.
\]

\subsection*{Задание 10}

Решите неравенство:
\[
4^x-5\cdot2^x+4>0.
\]

\newpage

\section*{Краткие решения тренировочного блока}

\subsection*{Задание 1}

\[
(3^x-3)(2-3^x)\ge0.
\]

Пусть
\[
t=3^x,\qquad t>0.
\]

Тогда:
\[
(t-3)(2-t)\ge0.
\]

Удобно вынести минус:
\[
-(t-3)(t-2)\ge0,
\]
\[
(t-3)(t-2)\le0.
\]

Отсюда:
\[
t\in[2;3].
\]

Возвращаемся к \(x\):
\[
2\le3^x\le3.
\]

Так как основание \(3>1\), получаем:
\[
\log_3 2\le x\le1.
\]

\[
\answer{\displaystyle x\in[\log_3 2;1].}
\]

\subsection*{Задание 2}

\[
\frac{4-2^x}{2^x-1}>0.
\]

Это неравенство разобрано выше.

\[
\answer{\displaystyle x\in(0;2).}
\]

\subsection*{Задание 3}

\[
\sqrt{x}\,(x-2)\le0.
\]

Область определения:
\[
x\ge0.
\]

На области определения:
\[
\sqrt{x}\ge0.
\]

Критические точки:
\[
\sqrt{x}=0 \Rightarrow x=0,
\]
\[
x-2=0 \Rightarrow x=2.
\]

Проверяем знаки на области \(x\ge0\):
\[
\begin{tikzpicture}[scale=1]
\draw[-{Stealth[length=2mm]}] (-0.5,0)--(5.5,0);

\draw (1,0.1)--(1,-0.1);
\draw (4,0.1)--(4,-0.1);
\node[below] at (1,-0.1) {\(0\)};
\node[below] at (4,-0.1) {\(2\)};

\draw[fill=black,draw=black,thick] (1,0) circle (2pt);
\draw[fill=black,draw=black,thick] (4,0) circle (2pt);

\node[above] at (2.5,0) {\(-\)};
\node[above] at (4.8,0) {\(+\)};

\draw[very thick,accent] (1,0.35)--(4,0.35);
\draw[fill=accent,draw=accent,thick] (1,0.35) circle (2pt);
\draw[fill=accent,draw=accent,thick] (4,0.35) circle (2pt);
\end{tikzpicture}
\]

\[
\answer{\displaystyle x\in[0;2].}
\]

\subsection*{Задание 4}

\[
5\cdot4^x-7\cdot10^x+2\cdot25^x<0.
\]

Это неравенство разобрано выше.

\[
\answer{\displaystyle x\in(0;1).}
\]

\subsection*{Задание 5}

\[
7\cdot3^{2x}+84\cdot3^x7^x-3\cdot7^{2x}<0.
\]

Делим на \(3^{2x}>0\):
\[
7+84\left(\frac73\right)^x-3\left(\frac73\right)^{2x}<0.
\]

Пусть
\[
t=\left(\frac73\right)^x,\qquad t>0.
\]

Тогда:
\[
7+84t-3t^2<0.
\]

Умножим на \(-1\) и поменяем знак:
\[
3t^2-84t-7>0.
\]

Найдём корни:
\[
D=(-84)^2-4\cdot3\cdot(-7)=7056+84=7140.
\]
\[
7140=4\cdot1785,
\qquad
\sqrt{7140}=2\sqrt{1785}.
\]

\[
t_{1,2}=\frac{84\pm2\sqrt{1785}}{6}
=
\frac{42\pm\sqrt{1785}}{3}.
\]

Так как
\[
\frac{42-\sqrt{1785}}{3}<0,
\]
а \(t>0\), отрицательный корень не подходит для положительной переменной \(t\).

Для неравенства
\[
3t^2-84t-7>0
\]
с учётом \(t>0\) получаем:
\[
t>\frac{42+\sqrt{1785}}{3}.
\]

Возвращаемся к \(x\):
\[
\left(\frac73\right)^x>\frac{42+\sqrt{1785}}{3}.
\]

Так как
\[
\frac73>1,
\]
получаем:
\[
x>\log_{\frac73}\frac{42+\sqrt{1785}}{3}.
\]

\[
\answer{\displaystyle x\in\left(\log_{\frac73}\frac{42+\sqrt{1785}}{3};+\infty\right).}
\]

\subsection*{Задание 6}

\[
25^x-6\cdot5^x+5\le0.
\]

Пусть
\[
t=5^x,\qquad t>0.
\]

Тогда:
\[
t^2-6t+5\le0.
\]

Разложим:
\[
t^2-6t+5=(t-1)(t-5).
\]

Получаем:
\[
(t-1)(t-5)\le0.
\]

Следовательно:
\[
t\in[1;5].
\]

Возвращаемся к \(x\):
\[
1\le5^x\le5.
\]

Так как
\[
1=5^0,\qquad 5=5^1,
\]
получаем:
\[
0\le x\le1.
\]

\[
\answer{\displaystyle x\in[0;1].}
\]

\subsection*{Задание 7}

\[
\frac{9^x-4}{3^x-1}\ge0.
\]

Пусть
\[
t=3^x,\qquad t>0.
\]

Тогда:
\[
9^x=(3^x)^2=t^2.
\]

Получаем:
\[
\frac{t^2-4}{t-1}\ge0.
\]

Разложим числитель:
\[
t^2-4=(t-2)(t+2).
\]

Так как \(t>0\), то:
\[
t+2>0.
\]

Значит, знак дроби определяется выражением:
\[
\frac{t-2}{t-1}\ge0.
\]

Критические точки:
\[
t=1 \quad \text{выколотая},
\]
\[
t=2 \quad \text{закрашенная}.
\]

С учётом \(t>0\):
\[
t\in(0;1)\cup[2;+\infty).
\]

Возвращаемся к \(x\):
\[
0<3^x<1
\]
или
\[
3^x\ge2.
\]

Так как основание \(3>1\), получаем:
\[
x<0
\]
или
\[
x\ge\log_3 2.
\]

\[
\answer{\displaystyle x\in(-\infty;0)\cup[\log_3 2;+\infty).}
\]

\subsection*{Задание 8}

\[
2^{x+1}+2^{x+2}-24<0.
\]

Преобразуем:
\[
2^{x+1}=2\cdot2^x,
\]
\[
2^{x+2}=4\cdot2^x.
\]

Тогда:
\[
2\cdot2^x+4\cdot2^x-24<0,
\]
\[
6\cdot2^x<24,
\]
\[
2^x<4.
\]

Так как
\[
4=2^2,
\]
получаем:
\[
x<2.
\]

\[
\answer{\displaystyle x\in(-\infty;2).}
\]

\subsection*{Задание 9}

\[
\left(\frac13\right)^x<9.
\]

Представим обе части как степени числа \(3\):
\[
\left(\frac13\right)^x=3^{-x},
\]
\[
9=3^2.
\]

Получаем:
\[
3^{-x}<3^2.
\]

Так как основание \(3>1\), сравниваем показатели:
\[
-x<2.
\]

Умножаем на \(-1\) и меняем знак:
\[
x>-2.
\]

\[
\answer{\displaystyle x\in(-2;+\infty).}
\]

\subsection*{Задание 10}

\[
4^x-5\cdot2^x+4>0.
\]

Пусть
\[
t=2^x,\qquad t>0.
\]

Тогда:
\[
4^x=(2^x)^2=t^2.
\]

Получаем:
\[
t^2-5t+4>0.
\]

Разложим:
\[
t^2-5t+4=(t-1)(t-4).
\]

Решаем:
\[
(t-1)(t-4)>0.
\]

С учётом \(t>0\):
\[
t\in(0;1)\cup(4;+\infty).
\]

Возвращаемся к \(x\):
\[
0<2^x<1
\]
или
\[
2^x>4.
\]

Значит:
\[
x<0
\]
или
\[
x>2.
\]

\[
\answer{\displaystyle x\in(-\infty;0)\cup(2;+\infty).}
\]

\newpage

\section*{Таблица ответов}

\begin{table}[h]
\caption{Ответы к тренировочному блоку}
\centering
\renewcommand{\arraystretch}{1.45}
\begin{tabularx}{\linewidth}{>{\centering\arraybackslash}p{16mm}X}
\toprule
\textbf{№} & \textbf{Ответ} \\
\midrule
1 & \(\displaystyle x\in[\log_3 2;1]\) \\
2 & \(\displaystyle x\in(0;2)\) \\
3 & \(\displaystyle x\in[0;2]\) \\
4 & \(\displaystyle x\in(0;1)\) \\
5 & \(\displaystyle x\in\left(\log_{\frac73}\frac{42+\sqrt{1785}}{3};+\infty\right)\) \\
6 & \(\displaystyle x\in[0;1]\) \\
7 & \(\displaystyle x\in(-\infty;0)\cup[\log_3 2;+\infty)\) \\
8 & \(\displaystyle x\in(-\infty;2)\) \\
9 & \(\displaystyle x\in(-2;+\infty)\) \\
10 & \(\displaystyle x\in(-\infty;0)\cup(2;+\infty)\) \\
\bottomrule
\end{tabularx}
\end{table}

\section*{Итоговая памятка}

\begin{enumerate}[label=\textbf{\arabic*.}, leftmargin=8mm]
\item В показательных неравенствах сначала ищем удобную замену.
\item После замены обязательно записываем условие \(t>0\).
\item Метод интервалов требует критических точек, знаков и выбора промежутков.
\item Точки из знаменателя всегда выколоты.
\item Корень чётной степени разрешает ноль: \(x\ge0\).
\item Если точка входит в ответ, нужна квадратная скобка.
\item Делить можно только на выражение, знак которого точно известен.
\item В системе нужно показывать пересечение множеств.
\item Если основание степени больше \(1\), знак неравенства между показателями сохраняется.
\item Если основание степени от \(0\) до \(1\), знак неравенства между показателями меняется.
\end{enumerate}

\end{document}